These dimensions can be approximated qualitatively or with lightweight heuristics; precision is less important than legibility for reflection.

These dimensions are not classifications to be assigned; they are lenses through which users can review their own sessions. A depth trace that drops suddenly after an error is not a ``low performance'' label---it is an invitation to ask: ``What happened there? Do I want to change something next time?''

\subsection{Reflective Artifacts, Not Measurements}

The distinction between measurement and reflective artifact is central to our proposal. A measurement aims for ground truth and produces verdicts (``your focus was low''). A reflective artifact aims for recognition and produces prompts (``your reasoning depth shifted here---does that match your experience?'').

Concretely, this means CPI systems using interaction traces should: (a) present patterns with uncertainty and alternative explanations, never single verdicts; (b) invite user annotation rather than inferring narratives; (c) support disagreement---a user who says ``that wasn't a breakdown, that was a strategic pause'' should be heard, not overridden; and (d) treat the user's interpretation as authoritative, not the system's.

\section{Design Vignettes}

We present two illustrative scenarios showing how interaction traces can complement wearable CPI data. These are design vignettes, not empirical findings---they are intended to make concrete the design space we are proposing.

\subsection{Vignette A: Workplace Deep-Work Session}

\textit{``The score says I failed. The trace says I recovered.''}

\textbf{Scenario.} A knowledge worker begins a morning deep-work session: drafting a technical proposal with an AI collaborator. Their wearable headband reports ``high stress'' and ``declining focus'' across the two-hour session. The aggregate score suggests declining focus across the session. The worker feels discouraged.

\textbf{With score alone.} The CPI dashboard shows a declining focus curve. The implied narrative is: the session went poorly, concentration eroded, the user should try shorter sessions or take more breaks. The user internalizes this as personal failure.

\textbf{With interaction traces added.} The conversation log tells a different story. For the first 45 minutes, reasoning depth was high and transitions were smooth---the user and AI were in sustained analytical flow. At 10:40am, an unexpected build failure triggered a rapid context switch: the user jumped to a debugging thread, breadth expanded as they pulled in three new knowledge domains, and depth dropped to quick diagnostic exchanges. By 11:15am, the user had resolved the issue and returned to the proposal---depth climbed back, breadth narrowed, and the session ended with a coherent draft.

\textbf{User annotation.} The user adds context: ``10:40 --- build broke, had to debug. Recovered well.'' The CPI system now shows: a session with a disruption and a successful recovery, not a session of declining focus.

\textbf{Reflection.} The user learns: my stress spikes during unexpected failures, but I recover within 35 minutes. Next time, I might schedule a 5-minute reset after disruptions rather than pushing through immediately. The score alone would have suggested ``try harder''; the trace suggests ``adjust the recovery pattern.''

\subsection{Vignette B: Creative Exploration Session}

\textit{``The system thinks I'm distracted. I'm actually incubating.''}

\textbf{Scenario.} A writer spends an afternoon exploring ideas for a new project with an AI collaborator. Their wearable reports ``low focus'' and ``high variability.'' A productivity-oriented CPI system nudges: ``You seem distracted. Would you like to start a focus session?'' The system has no notion of the user's session goal.

\textbf{With score alone.} The CPI system interprets low sustained focus as a problem. The nudge implies that the user should return to a ``productive'' state---sustained, narrow, deep. The user feels guilty about their exploratory afternoon.

\textbf{With interaction traces added.} The conversation log reveals a different cognitive rhythm: breadth is high and intentional---the user is deliberately connecting ideas across domains (literature, architecture, psychology). Transitions are frequent but rhythmic, not chaotic. And at three specific moments, depth spikes sharply: the user converges on an insight, explores it intensely for 10--15 minutes, then resumes broad exploration. This is not distraction; it is creative incubation with periodic convergence.

\textbf{User annotation.} The user marks the session as ``creative exploration'' and flags the three convergence moments: ``These were the breakthroughs.'' The CPI system learns: for this user, high-breadth sessions with punctuated depth spikes are productive, not distracted.

\textbf{Reflection.} The user adjusts their CPI settings: ``Creative sessions should not trigger focus nudges.'' The system updates the user's individualized baseline: exploration sessions have a different rhythm from deep-work sessions, and both are valid. The score would have prescribed one rhythm; the trace reveals that cognitive health requires multiple rhythms.

\section{Design Principles for Ethical CPI}

Drawing from the vignettes and broader considerations of user agency, we propose four design principles for CPI systems that integrate interaction traces.

\textbf{Principle 1: Uncertainty-first display.} CPI outputs should present confidence levels and alternative explanations by default, never single verdicts. \textit{Rationale:} Cognitive data is inherently ambiguous; false certainty erodes trust and agency. \textit{Example:} Instead of ``Focus: Low,'' display ``Focus estimate: low (confidence: moderate). Possible explanations: fatigue, creative exploration, task transition. What was happening?''

\textbf{Principle 2: User-contestable interpretations.} Any system-generated interpretation should include a visible mechanism for the user to disagree, reframe, or provide an alternative explanation---without penalty or loss of data. \textit{Rationale:} The user's lived experience is epistemically privileged over algorithmic inference. \textit{Example:} A ``This doesn't match my experience'' button that records the user's reinterpretation and adjusts future suggestions. The interface should preserve a trace of both the system's original suggestion and the user's reframe, creating an epistemic record of the interpretive dialogue.

\textbf{Principle 3: Annotation as first-class input.} User-provided context (reasons, events, intentions, emotional state) should be treated as primary data, not optional metadata. \textit{Rationale:} Interaction traces provide behavioral context; user annotations provide intentional context. Together, they create a richer reflective artifact than either alone. \textit{Example:} After a session, prompt ``Anything notable happen during this session?'' and integrate the respon